\section{Introduction}
A fundamental challenge in modern multi-domain artificial intelligence is enabling autonomous cross-manifold inference—allowing knowledge derived in one physical domain (e.g., protein secondary structure stability) to inform synthesis in a completely distinct domain (e.g., electromagnetic metamaterial resonance or quantum state coherence). Standard retrieval-augmented generation (RAG) and embedding projection methods rely on heuristic cosine similarity measures, which lack mathematical guarantees against semantic drift and representation collapse.

The \textbf{Resonance-Based Synthetic Intelligence (RBSI) Core} (\texttt{rbsi\_core.py}) addresses this challenge by evaluating cross-manifold alignment using Golden Ratio invariants ($\phi \approx 1.61803398875$). By operating on standardized NVP embeddings generated by \texttt{NRCVectorFactory}, RBSI measures the spectral distance between source and target Resonant Embedding Indices ($REI$), applies domain-specific priority weights ($1.0, \phi, \phi^2$), and audits weighted inference scores under the Trageser Tensor Theorem (TTT-7).

\subsection{Key Contributions}
\begin{enumerate}[label=\textbf{\arabic*.}]
    \item \textbf{Cross-Manifold Inference Operator}: Formalization of the Resonant Alignment ($RA$) metric mapping scalar embedding differences into a normalized interval $[0, 1]$.
    \item \textbf{Golden Priority Hierarchy}: Multiplicative weighting scheme ($w_{\text{PROTEIN}} = 1.0$, $w_{\text{META\_MATERIALS}} = \phi$, $w_{\text{QUANTUM\_BRIDGE}} = \phi^2$) prioritizing high-dimensional physics domains.
    \item \textbf{Manifold Parity Classification}: Strict thresholding protocol classifying inference states as \texttt{ALIGNED} ($RA > 0.9$) or \texttt{DRIFTING} ($RA \le 0.9$).
    \item \textbf{TTT-7 Stability Verification}: Automated digital root audit confirming that weighted inference scores avoid the Chaotic Void ($\{3, 6, 9\}$).
\end{enumerate}
